%!TEX root = ../QuantumNoiseAndDecoherence.tex
% ──────────────────────────────────────────────────────────────
%  Lecture 4 — Continuous Variables and Characteristic Functions
% ──────────────────────────────────────────────────────────────

% Continues Section 3 briefly, then opens Section 4.

\subsection{Single-mode Fock space and the harmonic oscillator}%
\label{sec:single-mode}

When the single-particle Hilbert space is $\hilb = \C$ (a single mode
of the electromagnetic field), Bose Fock space simplifies dramatically.
Since $\C^{\tp n} \cong \C$ for every~$n$ and the symmetrisation
projector acts as the identity, we obtain
\begin{equation}\label{eq:single-mode-fock-2}
  \mathcal{F}_+(\C)
  = \bigoplus_{n=0}^{\infty} \C
  \cong \ell^2(\N)\,.
\end{equation}
This is the Hilbert space of a single \emph{quantum harmonic
oscillator}, spanned by \textbf{Fock states} (number states)
$\ket{0}, \ket{1}, \ket{2}, \ldots$ The creation and annihilation
operators reduce to a single pair $a^\adj$, $a$ satisfying
\begin{equation}\label{eq:single-mode-ccr}
  [a,\, a^\adj] = \One\,,
\end{equation}
with the standard relations
\begin{equation}\label{eq:fock-state-relations}
  \ket{n} = \frac{(a^\adj)^n}{\sqrt{n!}}\ket{0}\,,\qquad
  a\ket{n} = \sqrt{n}\,\ket{n-1}\,,\qquad
  a^\adj\ket{n} = \sqrt{n+1}\,\ket{n+1}\,,\qquad
  a^\adj a\,\ket{n} = n\ket{n}\,.
\end{equation}

\begin{definition}[Continuous quantum variable]\label{def:cqv}
  A \emph{continuous quantum variable} (or \emph{mode}) is a quantum
  system whose Hilbert space is the single-mode Fock space
  $\mathcal{F}_+(\C) \cong \ell^2(\N)$.
\end{definition}

This is the relevant description for a single mode of light in an
optical cavity, a single phonon mode, or any quantum harmonic
oscillator.  In this course, both our system and its environment will
typically be modelled as collections of continuous quantum variables.

% ══════════════════════════════════════════════════════════════
\section{Continuous variables}\label{sec:continuous-variables}

We now develop the mathematical toolkit for working with continuous
quantum variables.  The key challenge: states and operators on the
infinite-dimensional space $\ell^2(\N)$ require infinitely many
parameters to specify.  We need compact, physically meaningful ways
to represent them.

% ──────────────────────────────────────────────────────────────
\subsection{Displacement operators}\label{sec:displacement}

\begin{definition}[Displacement operator]\label{def:displacement}
  For $\alpha \in \C$, the \textbf{displacement operator} is
  \begin{equation}\label{eq:displacement}
    \eqbox{D(\alpha) = \exp\!\bigl(\alpha\, a^\adj - \alpha^*\, a\bigr)}\,.
  \end{equation}
\end{definition}

Since the exponent is $i$ times a hermitian operator, $D(\alpha)$ is
\emph{unitary}: $D(\alpha)^\adj = D(\alpha)^{-1} = D(-\alpha)$.

\begin{intuition}[Why displacement operators?]
  Two reasons to prefer $D(\alpha)$ over the bare operators $a$,
  $a^\adj$:
  \begin{itemize}
    \item \textbf{Mathematical:} $a$ and $a^\adj$ are unbounded; $D(\alpha)$
      is bounded (unitary).  We can recover $a$ by differentiating
      $D(\alpha)$ with respect to~$\alpha$ at $\alpha = 0$.
    \item \textbf{Physical:} displacement operators correspond to
      linear optics transformations (beam splitters, phase shifters).
      You can buy a device that implements $D(\alpha)$; you cannot
      buy a device that implements~$a$.
  \end{itemize}
\end{intuition}

\begin{exercise}\label{ex:bch-displacement}
  Using the Baker--Campbell--Hausdorff formula, show that
  \begin{equation}\label{eq:displacement-bch}
    D(\alpha)
    = e^{-\lvert\alpha\rvert^2/2}\, e^{\alpha\, a^\adj}\,
      e^{-\alpha^*\, a}\,.
  \end{equation}
\end{exercise}

The displacement operators satisfy the \textbf{composition rule}:
\begin{equation}\label{eq:displacement-compose}
  D(\alpha)\,D(\beta)
  = e^{(\alpha\beta^* - \alpha^*\beta)/2}\,D(\alpha + \beta)\,.
\end{equation}

% ──────────────────────────────────────────────────────────────
\subsection{Coherent states}\label{sec:coherent-states}

\begin{definition}[Coherent state]\label{def:coherent-state}
  The \textbf{coherent state} $\ket{\alpha}$ is obtained by displacing
  the vacuum:
  \begin{equation}\label{eq:coherent-state}
    \ket{\alpha} = D(\alpha)\ket{0}
    = e^{-\lvert\alpha\rvert^2/2}
      \sum_{n=0}^{\infty} \frac{\alpha^n}{\sqrt{n!}}\,\ket{n}\,.
  \end{equation}
\end{definition}

Coherent states are the states produced by lasers and are central to
quantum optics.  They are eigenstates of the annihilation operator:
$a\ket{\alpha} = \alpha\ket{\alpha}$.

\subsubsection{Overcompleteness}

Coherent states form an \emph{overcomplete} set: they span the Hilbert
space, but are not orthogonal.  The \textbf{resolution of the identity}
reads
\begin{equation}\label{eq:coherent-completeness}
  \eqbox{\frac{1}{\pi}\int_{\C}
    \ketbra{\alpha}{\alpha}\;\dd^2\alpha = \One}\,,
  \qquad \dd^2\alpha = \dd(\mathrm{Re}\,\alpha)\;\dd(\mathrm{Im}\,\alpha)\,.
\end{equation}

Using this, any state $\ket{\psi}$ can be expanded as
\[
  \ket{\psi}
  = \frac{1}{\pi}\int_{\C} f(\alpha)\,\ket{\alpha}\;\dd^2\alpha\,,
  \qquad f(\alpha) = \braket{\alpha}{\psi}\,.
\]
Knowing the function $f(\alpha)$ is equivalent to knowing the state.
This is the simplest example of a \emph{quasi-distribution} function.

\begin{remark}
  Because the coherent states are overcomplete, the representation of
  a state in terms of $f(\alpha)$ is \emph{not unique}---different
  functions can yield the same state.  This redundancy is both a
  blessing (one can choose a convenient representative) and a curse
  (proving equality of two representations requires care).
\end{remark}

% ──────────────────────────────────────────────────────────────
\subsection{Characteristic functions}\label{sec:characteristic}

For \emph{mixed} states (density operators on an infinite-dimensional
space), we need a way to specify operators---not just vectors---in
terms of functions on~$\C$.

\subsubsection{Orthogonality of displacement operators}

The displacement operators satisfy the trace orthogonality relation
\begin{equation}\label{eq:displacement-orthogonality}
  \tr\bigl(D(\alpha)^\adj\, D(\beta)\bigr)
  = \pi\,\delta^{(2)}(\alpha - \beta)\,,
\end{equation}
where $\delta^{(2)}$ is the two-dimensional Dirac delta.  As a
consequence, any bounded operator~$A$ can be expanded as
\begin{equation}\label{eq:operator-expansion}
  \eqbox{A = \frac{1}{\pi}\int_{\C}
    \tr\bigl(A\, D(\alpha)^\adj\bigr)\, D(\alpha)\;\dd^2\alpha}\,.
\end{equation}

\subsubsection{The characteristic function}

\begin{definition}[Characteristic function]\label{def:char-fn}
  The \textbf{characteristic function} of a density operator~$\rho$ is
  \begin{equation}\label{eq:char-fn}
    \eqbox{\chi_\rho(\alpha) = \tr\bigl(\rho\, D(\alpha)\bigr)}\,,
    \qquad \alpha \in \C\,.
  \end{equation}
\end{definition}

\begin{keyresult}[Properties of the characteristic function]
  \leavevmode
  \begin{enumerate}[(i)]
    \item \textbf{Normalisation:} $\chi_\rho(0) = \tr(\rho) = 1$.
    \item \textbf{Hermiticity constraint:}
      $\chi_\rho(\alpha)^* = \chi_\rho(-\alpha)$.
    \item \textbf{Determines $\rho$:} the density operator can be
      reconstructed from its characteristic function via
      \begin{equation}\label{eq:rho-from-chi}
        \rho = \frac{1}{\pi}\int_{\C}
          \chi_\rho(\alpha)\, D(-\alpha)\;\dd^2\alpha\,.
      \end{equation}
    \item \textbf{Experimentally accessible:} $\chi_\rho(\alpha)$ can
      be estimated from homodyne and heterodyne detection data.
    \item \textbf{Related to the Wigner function:} writing
      $\alpha = (q + ip)/\sqrt{2}$, the Fourier transform of
      $\chi_\rho$ yields the Wigner function (see next lecture).
  \end{enumerate}
\end{keyresult}

% ──────────────────────────────────────────────────────────────
\subsection{Multiple modes}\label{sec:multiple-modes}

For $n$~distinguishable modes, the Hilbert space is
\begin{equation}\label{eq:multi-mode-hilb}
  \mathcal{F}_+(\C)^{\tp n}
  \cong \ell^2(\N)^{\tp n}\,,
\end{equation}
with independent creation and annihilation operators $a_j^\adj$, $a_j$
for $j = 1,\ldots,n$ satisfying $[a_j, a_k^\adj] = \delta_{jk}\One$.

The multi-mode displacement operator and coherent state are simply
tensor products:
\begin{align}
  D(\vb{\alpha})
  &= D(\alpha_1) \tp \cdots \tp D(\alpha_n)\,,
  &\vb{\alpha} &= (\alpha_1,\ldots,\alpha_n) \in \C^n\,,
  \label{eq:multi-mode-disp}\\
  \ket{\vb{\alpha}}
  &= \ket{\alpha_1} \tp \cdots \tp \ket{\alpha_n}\,.
  \label{eq:multi-mode-coherent}
\end{align}
The overcompleteness relation generalises to
\begin{equation}\label{eq:multi-mode-completeness}
  \frac{1}{\pi^n}\int_{\C^n}
    \ketbra{\vb{\alpha}}{\vb{\alpha}}\;\dd^{2n}\vb{\alpha} = \One\,,
\end{equation}
and the characteristic function of an $n$-mode state~$\rho$ is
$\chi_\rho(\vb{\alpha}) = \tr(\rho\, D(\vb{\alpha}))$.

\begin{remark}
  Multi-mode coherent states are always product states.  The richer
  class of \emph{Gaussian states}---which includes squeezed and
  entangled states---will be the subject of the next lecture.
\end{remark}
